%% ====================================================================
%%  atlas1.tex -- Optimization Atlas system overview (TikZ picture only).
%%  Compiled by atlas1-standalone.tex; \input directly into the paper.
%%
%%  ------------------------------------------------------------------
%%  LAYOUT SYSTEM  (read this before moving anything)
%%  ------------------------------------------------------------------
%%  The figure is four side-by-side PANELS reading left to right:
%%
%%    P1  Input / AtlasQuery      x  0.00 -- 3.10   centre 1.55
%%    P2  Parser & Seeds          x  3.65 -- 6.75   centre 5.20
%%    P3  SkyDiscover loop        x  8.05 -- 16.65  centre 12.35
%%    P4  Output / SolverRecipe   x 17.20 -- 20.30  centre 18.75
%%
%%  Inter-panel gaps are 0.55 except P2->P3, which is 1.30: that is the
%%  only inter-panel arrow carrying a text label ("initial population"),
%%  and the label has to sit clear of both panel borders.
%%
%%  All four panels share ONE vertical grid, so titles, subtitles and
%%  card rows line up straight across the whole figure:
%%
%%    \PanelTop  9.35   panel rectangle top
%%    \TitleY    8.88   panel title row               (all four panels)
%%    \SubY      8.28   panel subtitle row            (P2 has none)
%%    \ContentT  7.91   top edge of the content block (reference edge)
%%    \RowA      7.08   card row 1  (P1, P2 parse card, P3 top row, P4)
%%    \RowB      5.18   card row 2  (P1, P4)
%%    \RowC      3.28   card row 3  (P1, P4)
%%    \RowD      1.38   card row 4  (P1, P4, and P2's footnote)
%%    \BandY     2.68   horizontal FLOW BAND: the evaluator/selector
%%                      centre line, and the height at which the low
%%                      inter-panel arrows cross
%%    \ContentB  0.55   bottom edge of the content block (reference edge)
%%    \PanelBot  0.25   panel rectangle bottom
%%
%%  Rows A--D are 1.66 tall on a 1.90 pitch (0.24 clear between cards),
%%  so P1 and P4 are the same stack and P3's top row lands on row A too.
%%
%%  SIZE CLASSES.  No card gets an ad-hoc width.  Each picks one of the
%%  styles below (sidecard / loopcard / poolcard / chip), which fix text
%%  width and minimum height, so every card in a column is identical by
%%  construction.  Text widths carry >= 0.07cm of slack over the widest
%%  line they must hold, measured at the final font size.
%%
%%  LINE BREAKS.  Hyphenation is off and every break inside a card is
%%  manual, so no box can silently grow to fit its text.
%%
%%  ARROW RULE.  An inter-panel arrow attaches to a specific card when a
%%  specific card is the producer or consumer, and to the panel border
%%  when the panel as a whole is.  Intra-panel arrows always run
%%  node-anchor to node-anchor, so they cannot overlap what they connect.
%%
%%  FONTS.  Everything is \sffamily at an absolute \fontsize, so the
%%  figure is independent of the surrounding document's base size.  Note
%%  that submission.tex loads "times", which remaps \sffamily from CM
%%  Sans to Helvetica; Helvetica runs within about 1.7% of these metrics,
%%  every text width above still has slack, so the layout survives being
%%  \input into the paper.  Re-check the widest lines if the type scale
%%  or the text changes.
%%
%%  WHITESPACE.  Every line outside the tikzpicture ends in "%".  A bare
%%  line end there is a space token in the enclosing horizontal box and
%%  pads the figure's bounding box (it was worth 3.2cm before this).
%% ====================================================================
\begingroup%
%% --------------------------------------------------------------------
%% TYPE SCALE.  One command per text level; nothing in the picture picks
%% a font locally.  Leadings are tight (~1.15x) so the multi-line detail
%% text inside a card stays compact.
%% --------------------------------------------------------------------
\newcommand{\AtlasTitleFont}{\fontsize{10.5}{12.0}\bfseries\selectfont}%    panel title
\newcommand{\AtlasSubFont}{\fontsize{8.4}{9.6}\bfseries\selectfont}%       panel subtitle
\newcommand{\AtlasGroupFont}{\fontsize{7.8}{9.0}\bfseries\selectfont}%     in-panel group label
\newcommand{\AtlasHeadFont}{\fontsize{8.0}{9.2}\bfseries\selectfont}%      card heading
\newcommand{\AtlasCopyFont}{\fontsize{6.8}{7.8}\mdseries\selectfont}%      card detail
\newcommand{\AtlasStageHeadFont}{\fontsize{6.6}{7.6}\bfseries\selectfont}% funnel stage heading
\newcommand{\AtlasStageCopyFont}{\fontsize{6.0}{7.0}\mdseries\selectfont}% funnel stage detail
\newcommand{\AtlasNoteFont}{\fontsize{6.2}{7.2}\mdseries\selectfont}%      free-standing note
\newcommand{\AtlasChipFont}{\fontsize{7.0}{8.0}\mdseries\selectfont}%      processor-chip label
%
%% Card body = bold heading paragraph + detail paragraph.  The two parts
%% MUST be separate paragraphs, each with its \par inside the font group.
%% TeX fixes a paragraph's interline glue from \baselineskip as it stands
%% at \par, so a heading and its detail cannot share one paragraph: the
%% detail would inherit the document's 12pt leading.  (That was the cause
%% of the loose vertical spacing inside the old cards.)
\newcommand{\AtlasBody}[2]{%
  \parindent=0pt\relax%
  {\AtlasHeadFont\strut#1\par}%
  \vspace{0.45mm}%
  {\AtlasCopyFont\strut#2\strut\par}}%
%% Heading-only card (no detail line).
\newcommand{\AtlasBodyOnly}[1]{%
  \parindent=0pt\relax%
  {\AtlasHeadFont\strut#1\strut\par}}%
%% Same two-paragraph construction at funnel-stage sizes.
\newcommand{\AtlasStageBody}[2]{%
  \parindent=0pt\relax%
  {\AtlasStageHeadFont\strut#1\par}%
  \vspace{0.30mm}%
  {\AtlasStageCopyFont\strut#2\strut\par}}%
%
%% --------------------------------------------------------------------
%% GRID CONSTANTS.  Every coordinate in the picture is written in terms
%% of these, so a row or a column can be moved from one place.
%% --------------------------------------------------------------------
\def\PanelTop{9.35}%   panel rectangle, top edge (all four panels)
\def\PanelBot{0.25}%   panel rectangle, bottom edge (all four panels)
\def\TitleY{8.88}%     panel title row
\def\SubY{8.28}%       panel subtitle row
\def\ContentT{7.91}%   reference only: content top = top edge of row A
\def\ContentB{0.55}%   reference only: content bottom = row D / evaluator foot
\def\RowA{7.08}%       card row 1 centre
\def\RowB{5.18}%       card row 2 centre
\def\RowC{3.28}%       card row 3 centre
\def\RowD{1.38}%       card row 4 centre
\def\BandY{2.68}%      flow band / evaluator + selector centre line
\def\PoneL{0.00}%      P1 left edge
\def\PoneR{3.10}%      P1 right edge
\def\ColIn{1.55}%      P1 card column centre
\def\PtwoL{3.65}%      P2 left edge
\def\PtwoR{6.75}%      P2 right edge
\def\ColSeed{5.20}%    P2 card / chip column centre
\def\PthrL{8.05}%      P3 left edge
\def\PthrR{16.65}%     P3 right edge
\def\ColLoop{12.35}%   P3 centre, for the panel title only
\def\ColLoopA{10.20}%  P3 left  column: Context Builder / Solution Selector
\def\ColLoopB{14.50}%  P3 right column: Solution Generator / Evaluator
\def\PfourL{17.20}%    P4 left edge
\def\PfourR{20.30}%    P4 right edge
\def\ColOut{18.75}%    P4 card column centre
%
\begin{tikzpicture}[
  x=1cm, y=1cm, font=\sffamily, line cap=round, line join=round,
  %% Hyphenation off everywhere; see LINE BREAKS above.
  every node/.append style={execute at begin node={\hyphenpenalty=10000\relax}},
  %% ---- containers ---------------------------------------------------
  panel/.style={draw=black!72, line width=0.9pt, rounded corners=3.2mm},
  card/.style={draw=black!68, line width=0.75pt, rounded corners=1.8mm,
               align=center, inner sep=1.5mm},
  %% ---- card size classes (node width = text width + 2 x inner sep) ---
  %% sidecard 2.70 wide: holds "Compute Budget" (2.29) on one line.
  sidecard/.style={card, text width=2.40cm, minimum height=1.66cm},
  %% loopcard 3.70 wide: holds "bandit compute allocation" (2.86).
  loopcard/.style={card, text width=3.40cm, minimum height=1.66cm},
  %% poolcard: the selector.  Same width as a loopcard but taller, so the
  %% left column does not read as empty beside the tall evaluator.
  poolcard/.style={loopcard, minimum height=2.50cm},
  %% chip 1.55 wide: holds "Condat--V\~u" (1.28) on one line.
  chip/.style={card, rounded corners=0.6mm, inner sep=0.65mm,
               text width=1.42cm, minimum height=0.50cm,
               fill=atlasChip, text=white, font=\sffamily\AtlasChipFont},
  stagebox/.style={draw=black!70, line width=0.75pt},
  %% stage text: the caption sitting inside one funnel trapezoid.  Text
  %% width is set per stage, since each trapezoid is a different width.
  stage text/.style={align=center},
  %% ---- text-only nodes (font set on the node, so leading is correct) -
  panel title/.style={align=center, font=\sffamily\AtlasTitleFont},
  panel subtitle/.style={align=center, font=\sffamily\AtlasSubFont},
  group label/.style={align=center, font=\sffamily\AtlasGroupFont},
  note/.style={align=center, font=\sffamily\AtlasNoteFont},
  %% ---- connectors ----------------------------------------------------
  flow/.style={-{Latex[length=2.4mm,width=1.7mm]}, draw=black!66,
               line width=0.9pt},
  feedback/.style={-{Latex[length=2.4mm,width=1.7mm]}, draw=red!68!black,
                   line width=1.0pt}
]

%% --------------------------------------------------------------------
%% PALETTE.  Flat, print-safe tints.  Panels take the palest fills so the
%% cards inside stay legible; one hue per semantic role.
%% --------------------------------------------------------------------
\definecolor{atlasLavender}{HTML}{E7E0F1}% P1 panel  (input)
\definecolor{atlasSeedBg}{HTML}{DCE6F7}%   P2 panel  (parser + seeds)
\definecolor{atlasLoopBg}{HTML}{EAF0FA}%   P3 panel  (orchestration loop)
\definecolor{atlasOutBg}{HTML}{E3EFD7}%    P4 panel  (output)
\definecolor{atlasBlue}{HTML}{CFE0F5}%     card: problem statement
\definecolor{atlasTan}{HTML}{E9D9B7}%      card: data / parsing
\definecolor{atlasPurple}{HTML}{D8C8E8}%   card: hardware
\definecolor{atlasMint}{HTML}{C7F0DF}%     card: budget
\definecolor{atlasChip}{HTML}{668B90}%     seed-genome processor chips
\definecolor{atlasYellow}{HTML}{F4DFA5}%   card: LLM-driven stages
\definecolor{atlasGreen}{HTML}{CBE6BA}%    card: selection
\definecolor{atlasGray}{HTML}{E5E9EA}%     evaluator housing
\definecolor{atlasAmber}{HTML}{F0B923}%    funnel stage 0 (the hard gate)

%% ====================================================================
%% P1  INPUT PANEL -- "AtlasQuery": the four fields of a query.
%%     One group: a column of four identical sidecards on rows A--D.
%% ====================================================================
\path[panel, fill=atlasLavender] (\PoneL,\PanelBot) rectangle (\PoneR,\PanelTop);
\node[panel title]    at (\ColIn,\TitleY) {Input};
\node[panel subtitle] at (\ColIn,\SubY)   {AtlasQuery};

% P1 row A -- the objective to be solved, in any accepted surface form.
\node[sidecard, fill=atlasBlue] (problem) at (\ColIn,\RowA)
  {\AtlasBody{Problem}{natural language /\\CVXPY /\\conic data}};
% P1 row B -- what the synthesised solver will be tuned against.
\node[sidecard, fill=atlasTan] (instances) at (\ColIn,\RowB)
  {\AtlasBody{Instance Distribution}{real datasets}};
% P1 row C -- deployment target (heading only, no detail line).
\node[sidecard, fill=atlasPurple] (hardware) at (\ColIn,\RowC)
  {\AtlasBodyOnly{Hardware Target}};
% P1 row D -- search budget (heading only, no detail line).
\node[sidecard, fill=atlasMint] (budget) at (\ColIn,\RowD)
  {\AtlasBodyOnly{Compute Budget}};

%% ====================================================================
%% P2  PARSER & SEEDS PANEL.
%%     Two stacked groups: (a) the parse card on row A, (b) the "Seed
%%     Genomes" group -- label, four chips, one explanatory note.
%%     This panel has no subtitle, so the \SubY row is left empty on
%%     purpose rather than the title being dropped down to fill it.
%% ====================================================================
\path[panel, fill=atlasSeedBg] (\PtwoL,\PanelBot) rectangle (\PtwoR,\PanelTop);
\node[panel title] at (\ColSeed,\TitleY) {Parser \& Seeds};

% P2 row A -- lowers any accepted input form to the composite problem IR.
% On row A so it lines up with the Problem card that feeds it.
\node[sidecard, fill=atlasTan!55] (parser) at (\ColSeed,\RowA)
  {\AtlasBody{Parse $\rightarrow$\\Problem IR}{composite form\\$f+g+h(Lx)$}};

% P2 group label -- heads the chip stack below it; y = 5.88 sits
% midway between the parse card's foot (6.25) and the first chip (5.45).
\node[group label] at (\ColSeed,5.88) {Seed Genomes};

% P2 seed-genome chips -- four identical stylised processors on a 0.90
% pitch (5.20, 4.30, 3.40, 2.50).  The chip style fixes the body and the
% two \foreach loops add the pin rows, so all four are identical by
% construction; pin offsets are derived from the 1.55 x 0.50 body.
\foreach \id/\chiplabel/\cy in {%
  pdhg/PDHG/5.20, drs/DRS/4.30, fbs/FBS/3.40, cv/Condat--V\~u/2.50}{
  \node[chip] (\id) at (\ColSeed,\cy) {\chiplabel};
  \foreach \dx in {-0.56,-0.28,0,0.28,0.56}{% top and bottom pins
    \draw[black!55, line width=0.65pt] ([xshift=\dx cm]\id.north) -- +(0,0.10);
    \draw[black!55, line width=0.65pt] ([xshift=\dx cm]\id.south) -- +(0,-0.10);
  }
  \foreach \dy in {-0.13,0,0.13}{% side pins
    \draw[black!55, line width=0.65pt] ([yshift=\dy cm]\id.east) -- +(0.10,0);
    \draw[black!55, line width=0.65pt] ([yshift=\dy cm]\id.west) -- +(-0.10,0);
  }
}

% P2 note -- caption for the chip group.  Placed on row D so it aligns
% with the bottom card of P1 and P4.  Breaks are manual: the widest line
% is "textbook solvers as typed" at 2.39 against a 2.70 text width.
\node[note, text width=2.70cm] at (\ColSeed,\RowD)
  {textbook solvers as typed\\genomes; one island per\\base scheme};

%% ====================================================================
%% P3  SKYDISCOVER PANEL -- the orchestration loop.
%%     Internally a 2x2 arrangement:
%%       top-left  Context Builder      top-right  Solution Generator
%%       bot-left  Solution Selector    bot-right  Evaluator (housing)
%%     The cycle runs clockwise: context -> generator -> evaluator ->
%%     selector -> context, plus one red rejection arrow returning from
%%     the evaluator to the generator.
%% ====================================================================
\path[panel, fill=atlasLoopBg] (\PthrL,\PanelBot) rectangle (\PthrR,\PanelTop);
\node[panel title]    at (\ColLoop,\TitleY) {SkyDiscover};
\node[panel subtitle] at (\ColLoop,\SubY)   {Orchestration Loop};

% P3 top-left, row A -- assembles the prompt for the generator.
\node[loopcard, fill=atlasYellow] (context) at (\ColLoopA,\RowA)
  {\AtlasBody{Context Builder}{mutation menu +\\past results}};
% P3 top-right, row A -- the LLM proposal step.
\node[loopcard, fill=atlasYellow] (generator) at (\ColLoopB,\RowA)
  {\AtlasBody{Solution Generator}{LLM proposes genome\\edits as JSON patches}};
% P3 bottom-left -- the population manager.  Centred on \BandY like the
% evaluator, so the evaluator -> selector arrow is exactly horizontal.
\node[poolcard, fill=atlasGreen] (selector) at (\ColLoopA,\BandY)
  {\AtlasBody{Solution Selector}{dedup + islands +\\bandit compute allocation}};

%% --------------------------------------------------------------------
%% P3 bottom-right -- EVALUATOR: a housing card holding a group label and
%% a three-stage funnel.  The housing runs \ContentB (0.55) up to 4.81,
%% i.e. it is centred on \BandY like the selector.
%%
%% Funnel geometry.  Four horizontal edges, and the half-width at each:
%%     y = 4.26  half 1.710  (funnel mouth, 0.14 clear of the housing)
%%     y = 2.56  half 1.514
%%     y = 1.51  half 1.393
%%     y = 0.71  half 1.301  (funnel throat, 0.16 clear of the housing)
%% Half-width falls at a constant 0.1155 per unit height, so the two
%% sides are straight lines through all three stages.  Stage heights
%% (1.70 / 1.05 / 0.80) follow their line counts (5 / 3 / 2); stage 0 is
%% given the surplus so that the 0/1 seam clears \BandY, where the
%% outgoing arrow leaves, rather than meeting it.  Each stage's text
%% width sits below the narrower of its two edges, so no label can
%% spill outside its trapezoid.
%% --------------------------------------------------------------------
\node[card, fill=atlasGray, minimum width=3.70cm, minimum height=4.26cm]
  (evaluator) at (\ColLoopB,\BandY) {};
\node[group label] at (\ColLoopB,4.53) {Evaluator};% 0.28 below the housing top

% Funnel edge coordinates, top to bottom.  Named so the three trapezoids
% share them and cannot drift apart.
\coordinate (f0l) at (12.790,4.26); \coordinate (f0r) at (16.210,4.26);
\coordinate (f1l) at (12.986,2.56); \coordinate (f1r) at (16.014,2.56);
\coordinate (f2l) at (13.107,1.51); \coordinate (f2r) at (15.893,1.51);
\coordinate (f3l) at (13.199,0.71); \coordinate (f3r) at (15.801,0.71);

% Stage 0 -- the hard typing gate; amber because it is the only stage
% that rejects outright.  Widest stage, so it carries the most text.
\path[stagebox, fill=atlasAmber] (f0l) -- (f0r) -- (f1r) -- (f1l) -- cycle;
\node[stage text, text width=2.77cm] at (\ColLoopB,3.41)
  {\AtlasStageBody{Stage 0:\\Type-Check Gate}%
                  {convergence certificate\\required; invalid\\genomes rejected}};
% Stages 1--2 -- empirical measurement on real instances.
\path[stagebox, fill=black!7] (f1l) -- (f1r) -- (f2r) -- (f2l) -- cycle;
\node[stage text, text width=2.56cm] at (\ColLoopB,2.035)
  {\AtlasStageBody{Stages 1--2}%
                  {solve real instances;\\measure time-to-tolerance}};
% Stage 3 -- the final generalisation check; narrowest stage.
\path[stagebox, fill=black!3] (f2l) -- (f2r) -- (f3r) -- (f3l) -- cycle;
\node[stage text, text width=2.36cm] at (\ColLoopB,1.11)
  {\AtlasStageBody{Stage 3}{held-out + stress}};

%% --------------------------------------------------------------------
%% P3 connectors.  All five are axis-aligned and anchored to node
%% borders.  The two vertical arrows share the 1.44-tall corridor between
%% row A and the evaluator; they sit at +-1.25 from the column centre.
%% The caption is right-aligned against the red arrow, not centred between
%% the two, so it cannot be read as labelling the black one.
%% --------------------------------------------------------------------
\draw[flow] (context.east) -- (generator.west);                     % top edge, ->
\draw[flow] ([xshift=-1.25cm]generator.south)
         -- ([xshift=-1.25cm]evaluator.north);                      % right edge, down
\draw[flow] (evaluator.west) -- (selector.east);                    % bottom edge, <-
\draw[flow] (selector.north) -- (selector.north |- context.south);  % left edge, up
\draw[feedback] ([xshift=1.25cm]evaluator.north)
             -- ([xshift=1.25cm]generator.south);                   % rejection return, up
\node[note, text=red!68!black, anchor=east] at (15.62,5.53)% 0.13 left of the red
  {rejection feedback};%                            arrow, corridor mid-height

%% ====================================================================
%% P4  OUTPUT PANEL -- "SolverRecipe": four artefacts on the same stack
%%     geometry as P1, so rows A--D line up straight across the figure.
%%     Cards are plain white so they read as emitted artefacts rather
%%     than as the typed input slots of P1.
%% ====================================================================
\path[panel, fill=atlasOutBg] (\PfourL,\PanelBot) rectangle (\PfourR,\PanelTop);
\node[panel title]    at (\ColOut,\TitleY) {Output};
\node[panel subtitle] at (\ColOut,\SubY)   {SolverRecipe};

% P4 row A -- the winning genome and how it was reached.
\node[sidecard, fill=white] (best) at (\ColOut,\RowA)
  {\AtlasBody{Best Genome}{derivation +\\proof tree}};
% P4 row B -- the guarantees attached to it.
\node[sidecard, fill=white] (certificates) at (\ColOut,\RowB)
  {\AtlasBody{Convergence Certificates}{typing + rates +\\runtime residuals}};
% P4 row C -- the executable artefact.
\node[sidecard, fill=white] (compiled) at (\ColOut,\RowC)
  {\AtlasBody{Compiled JAX}{solver code}};
% P4 row D -- the reported evidence.
\node[sidecard, fill=white] (benchmark) at (\ColOut,\RowD)
  {\AtlasBody{Benchmark Card}{held-out results}};

%% ====================================================================
%% INTER-PANEL FLOW.  Source is the panel border where the whole panel
%% is the producer, and a card border where one card is.
%% ====================================================================
% P1 -> P2: the complete AtlasQuery enters the parser.  Runs along row A.
\draw[flow] (\PoneR,\RowA) -- (parser.west);
% P2 -> P3: the seed genome library becomes the loop's initial
% population.  Runs along the flow band; the label sits in the widened
% inter-panel gap, clear of both panel borders.
\draw[flow] (\PtwoR,\BandY) -- (selector.west);
\node[note, text width=1.10cm] at (7.40,3.04)% gap centre, 0.11 above the arrow
  {initial\\population};
% P3 -> P4: genomes that clear stage 3 are emitted as the SolverRecipe.
\draw[flow] (evaluator.east) -- (\PfourL,\BandY);

\end{tikzpicture}%
\endgroup%
